% Chapter 3: Methodology — CascadeLP architecture, tier routing, feature pipelines
\selectlanguage{greek}


\chapter{Μεθοδολογία}
\label{ch:methodology}


\section{Επισκόπηση Πλαισίου}
\label{sec:framework-overview}

Το \en{CascadeLP} είναι ένα τετραεπίπεδο πλαίσιο συμπερασμού βασισμένο σε αξιοπιστία
για πρόβλεψη ακμών σε γραφήματα με κειμενικά χαρακτηριστικά κόμβων.
Η κεντρική αρχή
σχεδιασμού είναι η \textit{προοδευτική κλιμάκωση}: ένα ζεύγος κόμβων αξιολογείται
πρώτα από τον υπολογιστικά φθηνότερο διαθέσιμο ταξινομητή, και μεταβιβάζεται στον
επόμενο ταξινομητή μόνο εάν η αξιοπιστία του πρώτου---μετρούμενη ως το μέγιστο των
δύο προβλεπόμενων πιθανοτήτων κλάσης---πέσει κάτω από ένα ρυθμισμένο όριο.
Στην παρούσα υλοποίηση όλα τα χαρακτηριστικά έχουν προϋπολογιστεί, επομένως η
δρομολόγηση μετρά ποια βαθμίδα λαμβάνει την απόφαση και όχι πόσες κωδικοποιήσεις
\en{Sentence-Transformer} εκτελέστηκαν.
Σε μια υλοποίηση κατ' απαίτηση, η ίδια
δρομολόγηση μπορεί να χρησιμοποιηθεί για να αποφεύγεται η ανάκτηση ή παραγωγή των
ακριβότερων χαρακτηριστικών για ζεύγη που έχουν ήδη επιλυθεί.

Τα τέσσερα επίπεδα του \en{CascadeLP} είναι τα εξής.
Το \textbf{Επίπεδο~0} είναι ένας
ντετερμινιστικός κανόνας που επιλύει αυτοβρόχους: οποιοδήποτε ζεύγος $(u, u)$ όπου
αμφότεροι οι τελικοί κόμβοι είναι ο ίδιος κόμβος αντιστοιχίζεται αμέσως στη θετική
ετικέτα.
Πρόκειται για κανόνα προσαρμοσμένο στο \en{DSAA 2023}, όπου όλοι εκτός
από έναν από τους επισημασμένους αυτοβρόχους είναι θετικοί· η μοναδική αρνητική
εγγραφή αντιμετωπίζεται ως ανωμαλία ετικέτας.
Αυτός ο κανόνας δεν απαιτεί
υπολογισμό χαρακτηριστικών ούτε συμπερασμό μοντέλου.
Το \textbf{Επίπεδο~1} είναι ο
\texttt{\en{StructuralClassifier}}: ένα μοντέλο Λογιστικής Παλινδρόμησης εκπαιδευμένο σε
τέσσερα δομικά ευρετικά χαρακτηριστικά (\en{Common Neighbours}, \en{Jaccard
Coefficient}, \en{Adamic-Adar Index} και \en{Preferential Attachment}). Είναι ταχύ,
ερμηνεύσιμο και αποτελεσματικό για ζεύγη με πλούσια επικάλυψη γειτονιάς.
Το \textbf{Επίπεδο~2} είναι ο \texttt{\en{PosClassifier}}: ένα μοντέλο Τυχαίου Δάσους
εκπαιδευμένο σε ένα 72-διάστατο διάνυσμα χαρακτηριστικών συχνότητας μερών του λόγου
που σχηματίζεται συνενώνοντας τα ιστογράμματα μερών του λόγου αμφότερων των τελικών
κόμβων.
Σχεδιάζεται για ζεύγη \en{cold-start} όπου τα δομικά χαρακτηριστικά
παρέχουν ελάχιστο ή καθόλου διακριτικό σήμα.
Το \textbf{Επίπεδο~3} είναι ο
\texttt{\en{EmbeddingClassifier}}: ένα μοντέλο Λογιστικής Παλινδρόμησης εκπαιδευμένο
στην ομοιότητα συνημίτονου μεταξύ των ενσωματώσεων \en{Sentence-Transformer} των
κειμενικών εγγράφων των δύο τελικών κόμβων.
Είναι το ακριβότερο επίπεδο ως προς την παραγωγή χαρακτηριστικών και λαμβάνει
απόφαση μόνο για ζεύγη που παραμένουν αβέβαια μετά τα Επίπεδα~1 και~2.
Στα
πειράματα οι ενσωματώσεις έχουν ήδη υπολογιστεί για όλους τους διαθέσιμους
κόμβους.

Σε κάθε επίπεδο, εάν η μέγιστη προβλεπόμενη πιθανότητα ισούται ή υπερβαίνει το
όριο $\tau_k$ του επιπέδου, το ζεύγος αντιστοιχίζεται στην ετικέτα με την υψηλότερη
πιθανότητα και η κλιμάκωση σταματά.
Εάν δεν ισχύει αυτό, το ζεύγος μεταβιβάζεται
στο επόμενο επίπεδο.
Τα ζεύγη που φτάνουν στο Επίπεδο~3 αντιστοιχίζονται πάντοτε
στην πρόβλεψη αυτού του επιπέδου· δεν υπάρχει περαιτέρω κλιμάκωση.
Σχηματικό
διάγραμμα της αλυσίδας φαίνεται στο Σχήμα~\ref{fig:cascade}.

\begin{figure}[h!]
    \centering
    \tikzset{
        block/.style={rectangle, draw, rounded corners, minimum width=3.6cm,
            minimum height=1cm, align=center},
        decision/.style={diamond, draw, aspect=2.6, align=center, inner sep=2pt},
        sink/.style={rectangle, draw, rounded corners, minimum width=2.6cm,
            minimum height=1cm, align=center, fill=black!5},
        escalate/.style={-{Stealth}, thick},
        exitpath/.style={-{Stealth}, thick, dashed},
        lanelabel/.style={font=\small, fill=white, inner sep=1pt},
    }
    \begin{tikzpicture}[node distance=1.1cm and 2.6cm]
        \node[block] (input) {Ζεύγος $(u,v)$};
        \node[decision, below=of input] (t0) {$u = v$;};
        \node[decision, below=of t0] (cs) {\en{cold-start};};
        \node[block, below=of cs] (t1) {Επίπεδο 1\\\en{StructuralClassifier}};
        \node[decision, below=of t1] (d1) {$\max(p_0,p_1)\!\geq\!\tau_1$;};
        \node[block, below=of d1] (t2) {Επίπεδο 2\\\en{PosClassifier}};
        \node[decision, below=of t2] (d2) {$\max(p_0,p_1)\!\geq\!\tau_2$;};
        \node[block, below=of d2] (t3) {Επίπεδο 3\\\en{EmbeddingClassifier}};
        \node[sink, right=5.6cm of t2] (sink) {Τελική\\ετικέτα};

        \draw[escalate] (input) -- (t0);
        \draw[escalate] (t0) -- node[left]{όχι} (cs);
        \draw[escalate] (cs) -- node[left]{όχι} (t1);
        \draw[escalate] (t1) -- (d1);
        \draw[escalate] (d1) -- node[left]{όχι} (t2);
        \draw[escalate] (t2) -- (d2);
        \draw[escalate] (d2) -- node[left]{όχι} (t3);

        % Four entry points on the sink's west edge, ordered top-to-bottom to
        % match the vertical order of their source nodes in the main chain.
        \coordinate (e1) at ($(sink.west)+(0,0.35)$); % from t0 (farthest above)
        \coordinate (e2) at ($(sink.west)+(0,0.12)$); % from d1 (adjacent above)
        \coordinate (e3) at ($(sink.west)+(0,-0.12)$); % from d2 (adjacent below)
        \coordinate (e4) at ($(sink.west)+(0,-0.35)$); % from t3 (farthest below)

        % Nested lanes: sources adjacent to the sink's row (d1, d2) bend soonest
        % (innermost); sources farther away (t0, t3) bend in an outer lane so
        % their vertical run never crosses an inner lane's horizontal run.
        \draw[exitpath] (d1.east) -- ++(1.2,0) |- (e2)
            node[lanelabel, pos=0.15]{ναι};
        \draw[exitpath] (d2.east) -- ++(1.2,0) |- (e3)
            node[lanelabel, pos=0.15]{ναι};
        \draw[escalate] (cs.east) -- ++(0.7,0) |- (t2.east)
            node[lanelabel, pos=0.08]{ναι};
        \draw[exitpath] (t0.east) -- ++(2.4,0) |- (e1)
            node[lanelabel, pos=0.2]{ναι, $\hat y=1$};
        \draw[exitpath] (t3.east) -- ++(2.4,0) |- (e4)
            node[lanelabel, pos=0.2]{πάντα};
    \end{tikzpicture}
    \caption{Σχηματική απεικόνιση της αλυσίδας \en{CascadeLP}: κάθε ζεύγος
    δρομολογείται μέσω διαδοχικών επιπέδων ταξινομητών, με ρητή παράκαμψη του
    δομικού επιπέδου για ζεύγη \en{cold-start} και πρόωρη έξοδο όταν η
    βαθμολογία εμπιστοσύνης υπερβαίνει το αντίστοιχο όριο.}
    \label{fig:cascade}
\end{figure}


\subsection{Αρχιτεκτονική Υλοποίησης}
\label{subsec:implementation-architecture}

Η περιγραφή του \en{CascadeLP} παραπάνω είναι αλγοριθμική· η υλοποίηση χωρίζεται σε
τέσσερα επίπεδα ευθύνης, καθένα εκφρασμένο ως ξεχωριστό πακέτο \en{Python}, ώστε η
εξαγωγή χαρακτηριστικών να μπορεί να τρέξει και να αποθηκευτεί ανεξάρτητα από την
εκπαίδευση ταξινομητών.

Η ροή ξεκινά από το \texttt{\en{src\slash{}data\slash{}loader.py}}.
Η συνάρτηση \texttt{\en{load\_edges()}} διαβάζει τα \texttt{\en{train.csv}}/\texttt{\en{test.csv}} απευθείας σε
μνήμη, αφού πρόκειται για μερικές εκατοντάδες χιλιάδες γραμμές· το \texttt{\en{nodes.tsv}}
όμως ζυγίζει 641\en{MB} και περιέχει το πλήρες κείμενο \en{Wikipedia} 837.855 άρθρων, οπότε
η \texttt{\en{load\_nodes\_for\_ids()}} το διαβάζει σε τμήματα (\en{chunks}) των 50.000 γραμμών
και κρατά μόνο τις σειρές των οποίων το \en{id} εμφανίζεται στο τρέχον σύνολο ζευγών —
έτσι αποφεύγεται η φόρτωση ολόκληρου του αρχείου στη μνήμη για κάθε \en{script} που
χρειάζεται μόνο ένα υποσύνολο κόμβων.
Η \texttt{\en{build\_graph()}} κατασκευάζει τον μη
κατευθυνόμενο γράφημα \en{NetworkX} από τα θετικά ζεύγη του συνόλου εκπαίδευσης· αυτό το
γράφημα τροφοδοτεί τόσο τα δομικά χαρακτηριστικά όσο και τη μάσκα \en{cold-start}.

Τρία αρχεία στο \texttt{\en{src\slash{}features\slash{}}} μετατρέπουν ζεύγη κόμβων σε αριθμητικά
διανύσματα.
Το \texttt{\en{structural.py}} υπολογίζει τους τέσσερις ευρετικούς δείκτες μέσω
της \texttt{\en{compute\_heuristics(G,\allowbreak{} pairs)}}, με το γράφημα πάντα ως πρώτο όρισμα.
Η αντιστροφή των ορισμάτων πρέπει να αποφεύγεται, επειδή η \en{NetworkX} μπορεί να
ερμηνεύσει το \en{DataFrame} ως πίνακα γειτνίασης και να οδηγήσει σε παραπλανητικό σφάλμα.
Το \texttt{\en{embeddings.py}} υλοποιεί
τον καθαρισμό κειμένου \en{Wikitext} (\texttt{\en{clean\_wiki\_text()}}) και τις δύο οικογένειες
σημασιολογικής ομοιότητας: \en{TF-IDF} μέσω \texttt{\en{build\_tfidf()}}/\texttt{\en{compute\_tfidf\_scores()}}, και
\en{Sentence-Transformer} μέσω \texttt{\en{encode\_nodes()}} (ενιαία κωδικοποίηση κάθε κόμβου, όχι
ανά ζεύγος) και \texttt{\en{compute\_embedding\_scores()}}.
Το \texttt{\en{linguistic.py}} παράγει το
72-διάστατο διάνυσμα συχνοτήτων μερών του λόγου μέσω \texttt{\en{compute\_pos\_features()}},
συνενώνοντας τα 36-διάστατα ιστογράμματα (\texttt{\en{pos\_frequency\_vector()}}) των δύο τελικών
κόμβων.
Τα \texttt{\en{scripts\slash{}data\slash{}compute\_structural.py}} και \texttt{\en{compute\_semantic.py}} είναι τα σημεία
εισόδου \en{CLI} που καλούν αυτές τις συναρτήσεις και αποθηκεύουν το αποτέλεσμα ως
\en{CSV} στο \texttt{\en{data\slash{}interim\slash{}}}, ώστε να υπολογίζονται μία φορά και να επαναχρησιμοποιούνται
σε κάθε πείραμα εκπαίδευσης.

Το \texttt{\en{src\slash{}models\slash{}svm.py}} περιέχει τους πέντε ταξινομητές του \en{sklearn}
(\texttt{\en{StructuralClassifier}}, \texttt{\en{TfidfClassifier}}, \texttt{\en{PosClassifier}},
\texttt{\en{EmbeddingClassifier}}, \texttt{\en{SvmClassifier}}), οι οποίοι μοιράζονται την ίδια διεπαφή
\texttt{\en{.fit()}}/\texttt{\en{.predict\_proba()}}/\texttt{\en{.save()}}/\texttt{\en{.load()}} — σχεδιαστική επιλογή που επιτρέπει
στο \texttt{\en{src\slash{}models\slash{}cascade.py}} να τους χειρίζεται εναλλάξιμα.
Η κλάση \texttt{\en{CascadeLP}}
κρατά τα Επίπεδα~1--3 ως τρία στιγμιότυπα αυτών των ταξινομητών· η
\texttt{\en{.fit()}} τα εκπαιδεύει το καθένα ανεξάρτητα στο πλήρες σύνολο εκπαίδευσης (χωρίς
αυτοβρόχους, που αφαιρούνται μέσω μάσκας), ενώ η \texttt{\en{.predict()}} επιστρέφει την τριάδα
\texttt{\en{(predictions,\allowbreak{} tier\_used,\allowbreak{} scores)}}.
Το δεύτερο διάνυσμα καταγράφει ποιο επίπεδο απάντησε σε
κάθε ζεύγος και το τρίτο την πιθανότητα της θετικής κλάσης από το επίπεδο που
έλαβε την απόφαση.
Αυτά καθιστούν δυνατή την ανάλυση δρομολόγησης του
Κεφαλαίου~\ref{ch:experiments}.
Η \texttt{\en{.tier\_stats()}} συνοψίζει τα ποσοστά δρομολόγησης ανά
επίπεδο, και η \texttt{\en{.save()}}/\texttt{\en{.load()}} σειριοποιούν ολόκληρο το αντικείμενο μέσω
\texttt{\en{joblib}} στο \texttt{\en{outputs\slash{}checkpoints\slash{}}}.

Για το αρχικό πείραμα \en{DSAA}, το \texttt{\en{scripts\slash{}analysis\slash{}run\_dsaa\_train.py}} συνδέει τα παραπάνω: φορτώνει τα
προϋπολογισμένα χαρακτηριστικά από το \texttt{\en{data\slash{}interim\slash{}dsaa\slash{}}}, διαχωρίζει σε υποσύνολα
εκπαίδευσης/επικύρωσης με στρωματοποίηση ως προς την ετικέτα
(\texttt{\en{train\_test\_split(...,\allowbreak{} stratify=y)}}), και εφαρμόζει το ίδιο κλείσιμο \texttt{\en{sub()}} σε
κάθε πίνακα χαρακτηριστικών ώστε όλα τα μοντέλα —συμπεριλαμβανομένου του
\en{CascadeLP}— να εκπαιδεύονται αποκλειστικά στους δείκτες \texttt{\en{tr}} και να αξιολογούνται
στους \texttt{\en{val}}.
Μετά την πρόβλεψη, το ίδιο \en{script} διασταυρώνει το
\texttt{\en{tier\_used}} με την ετικέτα δυσκολίας από το \texttt{\en{src\slash{}utils\slash{}difficulty.py}}
(\texttt{\en{label\_difficulty()}}) μέσω της \texttt{\en{tier\_difficulty\_breakdown()}} του
\texttt{\en{src\slash{}utils\slash{}metrics.py}}, και εξάγει ανά-ζεύγος στήλες
\texttt{\en{(tier\_used,\allowbreak{} difficulty,\allowbreak{} y\_pred)}} στο \texttt{\en{outputs\slash{}predictions\slash{}dsaa\slash{}cascade\_val\_tiers.csv}}.
Το
\texttt{\en{scripts\slash{}analysis\slash{}run\_dsaa\_evaluate.py}} επαναλαμβάνει την ίδια διασταύρωση στο σύνολο ελέγχου,
παράγοντας το αντίστοιχο \texttt{\en{cascade\_test\_tiers.csv}}.
Οι αναλύσεις ανά
επίπεδο και δυσκολία του \en{DSAA} προέρχονται από αυτά τα αρχεία.
Για το νέο
σύνολο, το \texttt{\en{scripts\slash{}analysis\slash{}run\_wiki\_cs\_8k\_experiment.py}}
χρησιμοποιεί την ομαδοποιημένη διαίρεση του \texttt{\en{src\slash{}data\slash{}protocol.py}},
κατασκευάζει γράφημα μόνο από το τμήμα εκπαίδευσης και αποθηκεύει τις μετρικές στο
\texttt{\en{outputs\slash{}stats\slash{}wiki\_cs\_8k\_experiment\_results.json}}.


\section{Πρωτόκολλο Ανίχνευσης Δεδομένων}
\label{sec:forensics-protocol}

Καμία αναφερόμενη βαθμολογία σε γράφημα με κειμενικά χαρακτηριστικά κόμβων δεν
πρέπει να γίνεται αποδεκτή χωρίς πρώτα να αποκλειστεί διαρροή δεδομένων και
τετριμμένη διαχωρισιμότητα, αφού και τα δύο μπορούν να προσομοιώσουν σημασιολογική
επίδοση που στην πραγματικότητα δεν υπάρχει.
Πριν από την εκπαίδευση οποιουδήποτε μοντέλου, εφαρμόζουμε επομένως ένα πρωτόκολλο
ανίχνευσης στο σύνολο δεδομένων \en{DSAA 2023} για να διακρίνουμε εάν οι σχεδόν
τέλειες βαθμολογίες που αναφέρονται στη βιβλιογραφία (Κεφάλαιο~\ref{ch:background},
\S\ref{sec:leakage-background}) οφείλονται σε διαρροή δεδομένων ή σε εγγενή
διαχωρισιμότητα του συνόλου δεδομένων.
Το πρωτόκολλο περιλαμβάνει τέσσερις ελέγχους.

\textbf{Έλεγχος ακριβούς επικάλυψης.} Κατασκευάζουμε το σύνολο των ζευγών $(u, v)$
του \texttt{\en{train.csv}} και ελέγχουμε πόσα ζεύγη του \texttt{\en{test.csv}}
εμφανίζονται ακριβώς σε αυτό το σύνολο.

\textbf{Έλεγχος αντεστραμμένης επικάλυψης.} Δεδομένου ότι το γράφημα αντιμετωπίζεται
ως μη κατευθυνόμενος (\S\ref{sec:level1-heuristics}), ελέγχουμε επιπλέον πόσα ζεύγη
$(u, v)$ του \texttt{\en{test.csv}} εμφανίζονται ως $(v, u)$ στο
\texttt{\en{train.csv}}.

\textbf{Ταυτοποίηση αυτοβρόχων.} Μετράμε τον αριθμό ζευγών όπου $u = v$ σε αμφότερα
τα αρχεία, και --- όπου υπάρχουν ετικέτες --- την κατανομή ετικετών αυτών των
ζευγών.
Ο κανόνας Επιπέδου~0 αντιμετωπίζει τον αυτοβρόχο ως θετικό, αλλά ο
έλεγχος καταγράφει και αποκλίσεις από αυτή τη σύμβαση, όπως τη μία αρνητική
εγγραφή του \en{DSAA 2023}.

\textbf{Ενδο-συνόλου διπλοεγγραφή.} Ελέγχουμε το \texttt{\en{train.csv}} για
διπλότυπα μη κατευθυνόμενα ζεύγη (το ίδιο $\{u, v\}$ σε παραπάνω από μία γραμμή),
αφού η τυχαία διαίρεση εκπαίδευσης-επαλήθευσης του \en{CascadeLP}
(\S\ref{sec:cascade-routing}) γίνεται βάσει γραμμής και όχι βάσει ταυτότητας
ζεύγους, οπότε διπλότυπα θα μπορούσαν να τοποθετήσουν το ίδιο φυσικό ζεύγος και
στις δύο πλευρές της διαίρεσης.

Πέραν του ελέγχου διαρροής, ορίζουμε ένα κριτήριο \textit{τετριμμένου ζεύγους} για
τον χαρακτηρισμό διαχωρισιμότητας: ένα ζεύγος χαρακτηρίζεται τετριμμένο εάν είναι
αυτοβρόχος, εάν ο αριθμός κοινών γειτόνων του υπερβαίνει ένα όριο $\theta_{\mathrm{CN}}$,
ή εάν η ομοιότητα συνημίτονου \en{TF-IDF} του υπερβαίνει ένα όριο
$\theta_{\mathrm{TF\text{-}IDF}}$.
Αντί για αυθαίρετα επιλεγμένα όρια, τα
$\theta_{\mathrm{CN}}$ και $\theta_{\mathrm{TF\text{-}IDF}}$ ορίζονται εμπειρικά ως
το εκατοστημόριο του πληθυσμού \textit{αρνητικών} ζευγών εκπαίδευσης πάνω από το
οποίο μόνο ένα κλάσμα $\alpha$ (προεπιλογή $\alpha = 0.01$) των αρνητικών παραμένει
— δηλαδή, το όριο πάνω από το οποίο η θετική/αρνητική κλάση γίνεται σχεδόν πλήρως
διαχωρίσιμη με βάση το χαρακτηριστικό αυτό μόνο του.
Ζεύγη που δεν πληρούν κανένα από τα τρία κριτήρια χαρακτηρίζονται λειτουργικά ως
«δύσκολα» (\en{hard}). Η ετικέτα δηλώνει μόνο αποτυχία των συγκεκριμένων απλών
κριτηρίων και δεν αποδεικνύει ότι απαιτείται σημασιολογική συλλογιστική.
Ο
Αλγόριθμος~\ref{alg:difficulty} περιγράφει τον υπολογισμό των ορίων και την
ανάθεση ετικέτας ανά ζεύγος· η προτεραιότητα των κριτηρίων (αυτοβρόχος $>$ υψηλό
\en{CN} $>$ υψηλή ομοιότητα κειμένου $>$ \en{hard}) αντικατοπτρίζει το κόστος της
κάθε οικογένειας χαρακτηριστικών, από τη φθηνότερη προς την ακριβότερη.

\begin{algorithm}[h!]
\caption{Πρωτόκολλο Χαρακτηρισμού Δυσκολίας}
\label{alg:difficulty}
\begin{algorithmic}[1]
\Require Ζεύγη εκπαίδευσης $D$ με ετικέτες $y$· \en{CN} και ομοιότητα \en{TF-IDF}
    για κάθε ζεύγος· ποσοστό $\alpha$ (προεπιλογή $0{,}01$)
\Ensure Ετικέτα δυσκολίας ανά ζεύγος
\State $\text{neg} \gets \{(u,v) \in D : y_{(u,v)} = 0\}$
\State $\theta_{\mathrm{CN}} \gets$ εκατοστημόριο $100(1-\alpha)$ του \en{CN} στο $\text{neg}$
\State $\theta_{\mathrm{TF\text{-}IDF}} \gets$ εκατοστημόριο $100(1-\alpha)$ της ομοιότητας \en{TF-IDF} στο $\text{neg}$
\For{$(u, v) \in D$}
    \If{$u = v$}
        \State ετικέτα $\gets$ \en{trivial\_self\_loop}
    \ElsIf{$\mathrm{CN}(u,v) > \theta_{\mathrm{CN}}$}
        \State ετικέτα $\gets$ \en{trivial\_high\_cn}
    \ElsIf{ομοιότητα \en{TF-IDF}$(u,v) > \theta_{\mathrm{TF\text{-}IDF}}$}
        \State ετικέτα $\gets$ \en{trivial\_high\_textsim}
    \Else
        \State ετικέτα $\gets$ \en{hard}
    \EndIf
\EndFor
\State \Return ετικέτες ανά ζεύγος
\end{algorithmic}
\end{algorithm}


\section{Προεπεξεργασία Δεδομένων}
\label{sec:preprocessing}

\subsection{Καθαρισμός Κειμένου}
\label{subsec:text-cleaning}

Και στα δύο σύνολα, η υλοποίηση δέχεται μία στήλη \texttt{\en{text}} ανά κόμβο.
Η συνάρτηση \texttt{\en{clean\_wiki\_text()}} αφαιρεί σχόλια \en{HTML}, στοιχεία
\texttt{\en{ref}}, λοιπές ετικέτες \en{HTML}, απλά πρότυπα \en{Wikitext} και
δείκτες έντονης ή πλάγιας γραφής.
Μετατρέπει τους εσωτερικούς συνδέσμους στη
λεζάντα τους, διατηρεί τη λεζάντα των εξωτερικών συνδέσμων και συμπτύσσει τα
διαδοχικά κενά.
Το ίδιο καθαρισμένο κείμενο χρησιμοποιείται από τα
χαρακτηριστικά \en{TF-IDF}, \en{POS} και \en{Sentence-Transformer}.

Η υλοποίηση δεν αποκωδικοποιεί ξεχωριστά οντότητες \en{HTML}, δεν αφαιρεί
συστηματικά οδηγούς προφοράς και δεν αναλύει δομημένο πεδίο \en{infobox}.
Επίσης
δεν κατασκευάζει εναλλακτικό κείμενο από τον τίτλο: όταν ένας κόμβος ή το κείμενό
του λείπει, τα χαρακτηριστικά \en{POS} είναι μηδενικά και οι ομοιότητες
\en{TF-IDF}/ενσωματώσεων λαμβάνουν την τιμή μηδέν.
Αυτή η συμπεριφορά αφορά
ιδιαίτερα τους \WfTextEmptyN{} κόμβους του νέου συνόλου και λαμβάνεται υπόψη ως
περιορισμός στην ερμηνεία των κειμενικών μοντέλων.

Ο κωδικοποιητής \en{Sentence-Transformer} εφαρμόζει εσωτερικά το δικό του μέγιστο
μήκος ακολουθίας.
Οι κανονικοποιημένες ενσωματώσεις υπολογίζονται μία φορά ανά
διαθέσιμο κόμβο και αποθηκεύονται, ενώ η ομοιότητα ενός ζεύγους προκύπτει από το
εσωτερικό γινόμενο των δύο διανυσμάτων.


\section{Επίπεδο 1 --- Δομικοί Ευρετικοί Δείκτες}
\label{sec:level1-heuristics}

\subsection{Ευρετικοί Δείκτες}
\label{subsec:heuristic-features}

Το Επίπεδο~1 αναπαριστά κάθε ζεύγος κόμβων $(u, v)$ ως τετραδιάστατο διάνυσμα
χαρακτηριστικών $\mathbf{f}_{\text{struct}}(u, v) = [\mathrm{CN}, \mathrm{Jaccard},
\mathrm{AA}, \mathrm{PA}]$ όπου οι τέσσερις συνιστώσες είναι ο αριθμός κοινών
γειτόνων (Εξίσωση~\ref{eq:cn}), ο συντελεστής \en{Jaccard} (Εξίσωση~\ref{eq:jaccard}),
ο δείκτης \en{Adamic-Adar} (Εξίσωση~\ref{eq:aa}), και η βαθμολογία \en{Preferential
Attachment} (Εξίσωση~\ref{eq:pa}).

Ο υπολογισμός αυτών των χαρακτηριστικών απαιτεί πρόσβαση στη γειτνιακή δομή του
γραφήματος.
Η ακριβής κατασκευή εξαρτάται από το πειραματικό πρωτόκολλο.
Στο αρχικό πείραμα
\en{DSAA}, κάθε θετικό ζεύγος του πλήρους \texttt{\en{train.csv}} συνεισφέρει
ακμή πριν από τη διαίρεση, επιλογή που δημιουργεί τον περιορισμό της
\S\ref{subsec:data-splits}.
Στο νέο σύνολο, το γράφημα κατασκευάζεται μόνο από
τις θετικές ακμές του τμήματος εκπαίδευσης και η εκάστοτε θετική ακμή-στόχος
αποκρύπτεται κατά τον υπολογισμό των χαρακτηριστικών της.
Η λίστα γειτνίασης αποθηκεύεται ως λεξικό που αντιστοιχίζει κάθε αναγνωριστικό
κόμβου στο σύνολο γειτόνων του, επιτρέποντας $O(|\Gamma(u)|)$ υπολογισμό της τομής
γειτονιάς για ένα ζεύγος ερωτήματος $(u, v)$.
Το γράφημα έχει \GraphNodes{} κόμβους με
τουλάχιστον έναν δομικό γείτονα.

Για κόμβους που δεν εμφανίζονται στο δομικό γράφημα (δηλαδή κόμβους χωρίς θετικά
ζεύγη εκπαίδευσης), τα χαρακτηριστικά \en{CN}, \en{Jaccard} και \en{AA} ορίζονται ως μηδέν, και το
χαρακτηριστικό \en{PA} είναι επίσης μηδέν αφού ο βαθμός του κόμβου είναι μηδέν.
Αυτό το
εκφυλισμένο διάνυσμα χαρακτηριστικών σηματοδοτεί στο Επίπεδο~1 ότι δεν διαθέτει
δομική πληροφορία.
Η υλοποίηση το αναγνωρίζει ρητά και παρακάμπτει το Επίπεδο~1,
δρομολογώντας το ζεύγος απευθείας στο Επίπεδο~2.
Ο Αλγόριθμος~\ref{alg:heuristics} περιγράφει
τον υπολογισμό των τεσσάρων χαρακτηριστικών για ένα σύνολο ζευγών.

\begin{algorithm}[h!]
\caption{Υπολογισμός Δομικών Ευρετικών Δεικτών}
\label{alg:heuristics}
\begin{algorithmic}[1]
\Require Γράφημα $G$ (μη κατευθυνόμενο, από θετικές ακμές εκπαίδευσης)· σύνολο
    ζευγών ερωτήματος $P$
\Ensure Πίνακας $\mathbf{f}_{\text{struct}}(u,v)$ για κάθε $(u,v) \in P$
\For{$(u, v) \in P$}
    \If{$u = v$}
        \State $\mathbf{f}_{\text{struct}}(u,v) \gets \text{NaN}$
            \Comment{επιλύεται ξεχωριστά στο Επίπεδο~0}
    \ElsIf{$u \notin G$ \textbf{ή} $v \notin G$}
        \State $\mathbf{f}_{\text{struct}}(u,v) \gets [0, 0, 0, 0]$
            \Comment{\en{cold-start}}
    \Else
        \State $\mathrm{CN} \gets |\Gamma(u) \cap \Gamma(v)|$
        \State $\mathrm{Jaccard} \gets |\Gamma(u) \cap \Gamma(v)| \,/\, |\Gamma(u) \cup \Gamma(v)|$
        \State $\mathrm{AA} \gets \sum_{w \,\in\, \Gamma(u) \cap \Gamma(v)} 1 / \log|\Gamma(w)|$
        \State $\mathrm{PA} \gets |\Gamma(u)| \cdot |\Gamma(v)|$
        \State $\mathbf{f}_{\text{struct}}(u,v) \gets [\mathrm{CN}, \mathrm{Jaccard}, \mathrm{AA}, \mathrm{PA}]$
    \EndIf
\EndFor
\State \Return $\mathbf{f}_{\text{struct}}$
\end{algorithmic}
\end{algorithm}

\subsection{Βαθμολόγηση Αξιοπιστίας}
\label{subsec:confidence-scoring}

Ο \en{StructuralClassifier} είναι ένα μοντέλο Λογιστικής Παλινδρόμησης με ποινή
κανονικοποίησης $L_2$, εκπαιδευμένο στα δομικά διανύσματα χαρακτηριστικών όλων των
ζευγών εκπαίδευσης.
Το μοντέλο παράγει ένα διάνυσμα πιθανότητας $[p_0, p_1]$ για
κάθε ζεύγος, όπου $p_0$ είναι η προβλεπόμενη πιθανότητα απουσίας ακμής και $p_1$
είναι η προβλεπόμενη πιθανότητα ύπαρξης ακμής.
Οι έξοδοι αυτές χρησιμοποιούνται ως βαθμολογίες εμπιστοσύνης για τη
δρομολόγηση.
Δεν εφαρμόστηκε ξεχωριστή βαθμονόμηση πιθανοτήτων ούτε αξιολόγηση με
διαγράμματα αξιοπιστίας, επομένως δεν ερμηνεύονται ως εγγυημένα βαθμονομημένες
πιθανότητες.

Η αξιοπιστία της πρόβλεψης του Επιπέδου~1 για ένα ζεύγος ορίζεται ως:
\begin{equation}
    c_1(u, v) = \max(p_0, p_1).
    \label{eq:confidence-tier1}
\end{equation}
Εάν $c_1(u, v) \geq \tau_1$, το ζεύγος αντιστοιχίζεται στην ετικέτα $\arg\max(p_0,
p_1)$ και η κλιμάκωση σταματά.
Διαφορετικά, το ζεύγος προωθείται στο Επίπεδο~2.
Το όριο $\tau_1$ ορίζεται χειρωνακτικά στην τιμή 0{,}8, μια συνηθισμένη επιλογή
συντηρητικής αξιοπιστίας για αλυσίδες κλιμάκωσης βασισμένες σε πιθανότητα.
Η
ευαισθησία της συνολικής απόδοσης ως προς αυτή την επιλογή εξετάζεται εκ των
υστέρων μέσω μελέτης ablation σε πλέγμα τιμών στο σύνολο επαλήθευσης
(\S\ref{sec:ablation}), όχι μέσω μιας αυτοματοποιημένης διαδικασίας επιλογής που
καθορίζει την τελική τιμή.


\section{Επίπεδο 2 --- Ελαφρύς Σημασιολογικός Ταξινομητής}
\label{sec:tier2}

\subsection{Χαρακτηριστικά Μερών του Λόγου + \en{Random Forest}}
\label{subsec:pos-features}

Ο προεπιλεγμένος ταξινομητής Επιπέδου~2 είναι ο \texttt{\en{PosClassifier}}, ένα \en{Random Forest}~\cite{breiman2001random}
εκπαιδευμένο σε 72-διάστατα διανύσματα χαρακτηριστικών συχνότητας μερών του
λόγου.
Για κάθε κόμβο $v$, το καθαρισμένο κειμενικό έγγραφο $d_v$ αποδίδεται σε
\en{tokens} με τον \en{tokenizer} \texttt{\en{punkt}} του \en{NLTK} και επισημαίνεται
με μέρη του λόγου χρησιμοποιώντας τον \texttt{\en{averaged\_perceptron\_tagger}} του
\en{NLTK}, που αντιστοιχίζει ετικέτα μέρους του λόγου του συνόλου \en{Penn Treebank}
(\en{NN, VB, JJ, RB, NNP} κ.λπ.) σε κάθε \en{token}.
Το διάνυσμα συχνότητας μερών του
λόγου $\mathbf{p}_v \in \mathbb{R}^{36}$ σχηματίζεται μετρώντας τον αριθμό \en{tokens}
που αντιστοιχίζονται σε καθεμία από τις 36 ετικέτες \en{Penn Treebank} που
χρησιμοποιούνται σε αυτή την εργασία και κανονικοποιώντας με τον συνολικό αριθμό
\en{tokens}, δίνοντας σχετικές συχνότητες.
Το διάνυσμα χαρακτηριστικών ζεύγους είναι
η συνένωση των δύο διανυσμάτων τελικών κόμβων: $\mathbf{f}_{\text{POS}}(u,
v) = [\mathbf{p}_u \,\|\, \mathbf{p}_v] \in \mathbb{R}^{72}$.
Η συνένωση διατηρεί τη σειρά των άκρων, παρότι το γράφημα αξιολογείται ως μη
κατευθυνόμενο.
Συνεπώς η πρόβλεψη μπορεί να αλλάξει αν ανταλλαχθούν τα $u$ και
$v$.
Δεν εφαρμόστηκε συμμετρικοποίηση μέσω μέσου όρου των δύο διατάξεων, και η
εξάρτηση αυτή αποτελεί περιορισμό της παρούσας υλοποίησης.

Ο \en{PosClassifier} είναι \en{Random Forest} με \PosNEstimators{} δέντρα και στάθμιση κλάσεων
\en{«balanced»} (αντίστροφη προς τη συχνότητα κάθε κλάσης, ώστε να αντισταθμίζεται η
ανισορροπία \TrainPosPct\%/\TrainNegPct\% θετικών/αρνητικών ζευγών), με τις λοιπές υπερπαραμέτρους στις
προεπιλεγμένες τιμές του \en{scikit-learn}~\cite{pedregosa2011scikit}.
Εκπαιδεύεται στα \TrainPairsNoSelf{} ζεύγη
εκπαίδευσης που παραμένουν μετά την εξαίρεση των \SelfLoopsTrain{} αυτοβρόχων
(\S\ref{sec:leakage-results}), οι οποίοι επιλύονται ξεχωριστά στο Επίπεδο~0.
Τα
Τυχαία Δάση είναι κατάλληλα
για αυτόν τον χώρο χαρακτηριστικών επειδή το 72-διάστατο διάνυσμα \en{POS} περιέχει
πολλά σχεδόν μηδενικά χαρακτηριστικά για ασυνήθεις κατηγορίες μερών του λόγου, και
η δενδροειδής επιλογή χαρακτηριστικών χειρίζεται φυσικά αυτή την αραιότητα.
Το
μοντέλο είναι επίσης ανθεκτικό στις διαφορές κλίμακας μεταξύ χαρακτηριστικών
συχνότητας μεγάλων άρθρων (που διαθέτουν καλά εκτιμημένες κατανομές \en{POS}) και
μικρών άρθρων (που ενδέχεται να έχουν θορυβώδεις κατανομές λόγω μικρού δείγματος).

Η βαθμολόγηση αξιοπιστίας για το Επίπεδο~2 ακολουθεί την ίδια σύμβαση με το
Επίπεδο~1: η αξιοπιστία $c_2(u, v) = \max(p_0, p_1)$ συγκρίνεται με το όριο $\tau_2$,
και τα ζεύγη κάτω από το όριο κλιμακώνονται στο Επίπεδο~3.


\section{Επίπεδο 3 --- Εκλέπτυνση Μέσω Μετασχηματιστή}
\label{sec:tier3}

\subsection{Προσαρμογή \en{Sentence-Transformer}}
\label{subsec:sentence-transformer}

Το Επίπεδο~3 είναι ο \texttt{\en{EmbeddingClassifier}}, που χρησιμοποιεί έναν κωδικοποιητή
\en{Sentence-Transformer} για την παραγωγή πυκνών σημασιολογικών ενσωματώσεων
κειμενικών εγγράφων κόμβων και στη συνέχεια εκπαιδεύει έναν ταξινομητή Λογιστικής
Παλινδρόμησης στην ομοιότητα συνημίτονου μεταξύ των δύο ενσωματώσεων τελικών κόμβων.

Το μοντέλο \en{Sentence-Transformer} που χρησιμοποιείται είναι το
\texttt{\en{\EmbeddingModelName{}}}, ένα ελαφρύ μοντέλο προσαρμοσμένο από αρχιτεκτονική
\en{MiniLM} σε μεγάλη συλλογή εργασιών ομοιότητας προτάσεων χρησιμοποιώντας
\en{contrastive learning}.
Παράγει 384-διάστατες ενσωματώσεις προτάσεων και εκτελείται
σημαντικά ταχύτερα από πλήρεις κωδικοποιητές \en{BERT} ή \en{RoBERTa} διατηρώντας
ισχυρή απόδοση ομοιότητας σημασιολογικής.
Για κάθε κόμβο $v$, η ενσωμάτωση
$\mathbf{e}_v = \mathrm{ST}(d_v) \in \mathbb{R}^{384}$ υπολογίζεται μία φορά και
αποθηκεύεται.
Κατά τον χρόνο αξιολόγησης ζεύγους, το χαρακτηριστικό που
χρησιμοποιείται είναι η βαθμωτή ομοιότητα συνημίτονου:
\begin{equation}
    \mathrm{sim}(u, v) = \frac{\mathbf{e}_u \cdot \mathbf{e}_v}
    {\|\mathbf{e}_u\| \cdot \|\mathbf{e}_v\|}.
    \label{eq:cosine-sim}
\end{equation}
Ο \en{EmbeddingClassifier} είναι ένα μοντέλο Λογιστικής Παλινδρόμησης εκπαιδευμένο
σε αυτό το ενιαίο χαρακτηριστικό.
Παρά την απλότητά του---ένα μονοδιάστατο
χαρακτηριστικό που τροφοδοτείται σε ένα γραμμικό μοντέλο---αυτός ο ταξινομητής
καταγράφει αποτελεσματικά τη σημασιολογική ευθυγράμμιση μεταξύ ζευγών κειμένου
κόμβων, επειδή ο χώρος ενσωμάτωσης του \en{Sentence-Transformer} εκπαιδεύεται ειδικά
ώστε να τοποθετεί σημασιολογικά παρόμοια κείμενα κοντά μεταξύ τους.
Η επιλογή
Λογιστικής Παλινδρόμησης αντί πιο σύνθετου μοντέλου (π.χ., \en{MLP} ή \en{SVM}) είναι
σκόπιμη: με ένα μόνο χαρακτηριστικό, ένα πιο σύνθετο μοντέλο θα υπερπροσαρμοζόταν,
και οι συνεχείς έξοδοι πιθανότητας της Λογιστικής Παλινδρόμησης παρέχουν τη
βαθμολογία της θετικής κλάσης.
Εφόσον το Επίπεδο~3 είναι τελικό, η έξοδος αυτή
δεν χρησιμοποιείται για περαιτέρω δρομολόγηση.

Το Επίπεδο~3 δεν έχει περαιτέρω κλιμάκωση· όλα τα ζεύγη που το φτάνουν λαμβάνουν
την πρόβλεψη του \en{EmbeddingClassifier} ανεξάρτητα από αξιοπιστία.


\section{Δρομολόγηση Βάσει Αξιοπιστίας}
\label{sec:cascade-routing}

Ο Αλγόριθμος~\ref{alg:cascade} τυποποιεί τη διαδικασία δρομολόγησης \en{CascadeLP}.

\begin{algorithm}[h!]
\caption{Δρομολόγηση \en{CascadeLP} για ζεύγος $(u, v)$}
\label{alg:cascade}
\begin{algorithmic}[1]
\Require Ζεύγος κόμβων $(u, v)$· όρια $\tau_1, \tau_2$
\Ensure Προβλεπόμενη ετικέτα $\hat{y} \in \{0, 1\}$
\If{$u = v$}
    \State \Return $\hat{y} \gets 1$ \Comment{Επίπεδο 0: αυτοβρόχος}
\EndIf
\State $\mathbf{f}_{\text{struct}} \gets$ υπολογισμός δομικών χαρακτηριστικών$(u, v)$
\If{$\mathbf{f}_{\text{struct}} \neq [0,0,0,0]$} \Comment{παράκαμψη στο \en{cold-start}}
    \State $(p_0, p_1) \gets \text{\en{StructuralClassifier}.predict\_proba}(\mathbf{f}_{\text{struct}})$
    \If{$\max(p_0, p_1) \geq \tau_1$}
        \State \Return $\hat{y} \gets \arg\max(p_0, p_1)$ \Comment{Επίπεδο 1: επίλυση}
    \EndIf
\EndIf
\State $\mathbf{f}_{\text{POS}} \gets$ υπολογισμός χαρακτηριστικών \en{POS}$(u, v)$
\State $(p_0, p_1) \gets \text{\en{PosClassifier}.predict\_proba}(\mathbf{f}_{\text{POS}})$
\If{$\max(p_0, p_1) \geq \tau_2$}
    \State \Return $\hat{y} \gets \arg\max(p_0, p_1)$ \Comment{Επίπεδο 2: επίλυση}
\EndIf
\State $s \gets \mathrm{sim}(\mathbf{e}_u, \mathbf{e}_v)$ \Comment{αναζήτηση προ-υπολογισμένων ενσωματώσεων}
\State \Return $\hat{y} \gets \text{\en{EmbeddingClassifier}.predict}(s)$ \Comment{Επίπεδο 3: πάντοτε επιλύει}
\end{algorithmic}
\end{algorithm}

Για κάθε ζεύγος $(u, v)$ στο σύνολο ελέγχου, ο αλγόριθμος προχωρά ως εξής.
Πρώτον,
εάν $u = v$, επιστρέφει αμέσως ετικέτα 1 (Επίπεδο~0).
Δεύτερον, υπολογίζει το δομικό διάνυσμα χαρακτηριστικών
$\mathbf{f}_{\text{struct}}(u, v)$.
Αν και τα τέσσερα στοιχεία είναι μηδενικά,
το ζεύγος αναγνωρίζεται ως \en{cold-start} και παρακάμπτει το Επίπεδο~1.
Διαφορετικά λαμβάνει το διάνυσμα προβλεπόμενης πιθανότητας του
\en{StructuralClassifier}.
Εάν $\max(p_0, p_1)
\geq \tau_1$, επιστρέφει την αντίστοιχη ετικέτα (επίλυση Επιπέδου~1).
Τρίτον,
υπολογίζει το διάνυσμα χαρακτηριστικών \en{POS} $\mathbf{f}_{\text{POS}}(u, v)$ και
λαμβάνει το διάνυσμα προβλεπόμενης πιθανότητας του \en{PosClassifier}.
Εάν $\max(p_0,
p_1) \geq \tau_2$, επιστρέφει την αντίστοιχη ετικέτα (επίλυση Επιπέδου~2).
Τέλος,
αναζητά τις προ-υπολογισμένες ενσωματώσεις \en{Sentence-Transformer} για τους $u$ και
$v$, υπολογίζει την ομοιότητά τους συνημίτονου, και επιστρέφει την προβλεπόμενη
ετικέτα του \en{EmbeddingClassifier} (επίλυση Επιπέδου~3).

Τα όρια $\tau_1 = 0{,}8$ και $\tau_2 = 0{,}7$ είναι χειρωνακτικά επιλεγμένες
προεπιλογές του \texttt{\en{CascadeLP}}, όχι το αποτέλεσμα αυτοματοποιημένης
βελτιστοποίησης.
Η ευαισθησία της συνολικής βαθμολογίας \en{Macro F1} και του
ρυθμού δρομολόγησης στο Επίπεδο~3 ως προς αυτή την επιλογή εξετάζεται εκ των υστέρων στο
σύνολο επαλήθευσης (20\% των ζευγών εκπαίδευσης, διατηρούμενο κατά στρωματοποίηση
ανά ετικέτα), σαρώνοντας ένα αραιό πλέγμα τιμών $\tau_1 \in \{0{,}6, 0{,}7, 0{,}8,
0{,}9, 0{,}95\}$ και $\tau_2 \in \{0{,}5, 0{,}6, 0{,}7, 0{,}8, 0{,}9\}$, χωρίς
επανεκπαίδευση των ταξινομητών ανά επίπεδο --- τα όρια επηρεάζουν μόνο τη
δρομολόγηση κατά το συμπερασμό, όχι την εκπαίδευση (\S\ref{sec:ablation}). Αυτό
χαρακτηρίζει την ανταλλαγή ακρίβειας-κόστους αντί να επιλέγει αυτόματα ένα
βέλτιστο σημείο λειτουργίας.

Μια σημαντική λεπτομέρεια υλοποίησης είναι ότι τα δομικά, \en{POS} και
ενσωματωτικά χαρακτηριστικά προϋπολογίζονται και αποθηκεύονται πριν από την
απόφαση δρομολόγησης.
Η μέτρηση χρόνου συμπερασμού περιλαμβάνει συνεπώς την
ανάγνωση αυτών των πινάκων και τις προβλέψεις των ταξινομητών, αλλά όχι την
εξαγωγή χαρακτηριστικών.
Τα ποσοστά ανά επίπεδο τεκμηριώνουν τον καταμερισμό των
αποφάσεων.
Η μετατροπή τους σε πραγματική εξοικονόμηση κωδικοποίησης απαιτεί
ξεχωριστή, τεμπέλικη υλοποίηση που θα παράγει τα ακριβότερα χαρακτηριστικά μόνο
μετά την κλιμάκωση.


\section{Χειρισμός \en{Cold-Start}}
\label{sec:cold-start}

Ένα \textit{ζεύγος \en{cold-start}} είναι ένα ζεύγος $(u, v)$ όπου
τουλάχιστον ένα από τα $u$ ή $v$ δεν έχει δομικούς γείτονες στο γράφημα
εκπαίδευσης---δηλαδή όπου $|\Gamma(u)| = 0$ ή $|\Gamma(v)| = 0$.
Για τέτοια ζεύγη,
τα χαρακτηριστικά \en{Common Neighbours}, \en{Jaccard} και \en{Adamic-Adar} είναι
ταυτολογικά μηδέν, και η βαθμολογία \en{Preferential Attachment} είναι επίσης μηδέν.
Η συνάρτηση \texttt{\en{cold\_start\_from\_structural()}} ανιχνεύει αυτό το
τετραπλό μηδενικό διάνυσμα.
Το \en{CascadeLP} παρακάμπτει ρητά τον
\en{StructuralClassifier} για αυτές τις γραμμές και τις εισάγει στο Επίπεδο~2
μαζί με όσα ζεύγη κλιμακώθηκαν λόγω χαμηλής εμπιστοσύνης στο Επίπεδο~1.

Τα χαρακτηριστικά \en{POS} του Επιπέδου~2 υπολογίζονται από το κείμενο κόμβου και
είναι επομένως διαθέσιμα για κόμβους \en{cold-start} ανεξάρτητα από την
κατάστασή τους στο δομικό γράφημα.
Αυτό είναι το κύριο πλεονέκτημα του πολυεπίπεδου
σχεδιασμού: το Επίπεδο~2 μπορεί να επιλύσει ζεύγη \en{cold-start} που το
Επίπεδο~1 δεν μπορεί, χρησιμοποιώντας μια οικογένεια χαρακτηριστικών που είναι
εντελώς ανεξάρτητη από την τοπολογία του γραφήματος.
Εάν τα χαρακτηριστικά \en{POS}
αποτύχουν επίσης να παράγουν σίγουρη πρόβλεψη---κάτι που συμβαίνει όταν αμφότεροι
οι κόμβοι έχουν πολύ σύντομα ή πανομοιότυπα κειμενικά έγγραφα---το ζεύγος κλιμακώνεται
στο Επίπεδο~3, όπου ο \en{Sentence-Transformer} παρέχει τη βαθύτερη σημασιολογική
σύγκριση.

Ένας πρόχειρος υπολογισμός στο σύνολο δεδομένων \en{DSAA 2023} δείχνει ότι οι κόμβοι
που απουσιάζουν εντελώς από το δομικό γράφημα ανέρχονται σε περίπου 170.000 από τους
συνολικά \GraphNodesTotal{} κόμβους (\GraphNodesTotal{} $-$ \GraphNodes{} $\approx$ 170.000)· πρόκειται για απλή
αφαίρεση, όχι για ανεξάρτητα μετρηθέν μέγεθος.
Η ανάλυση της \S\ref{sec:cold-start-results} χρησιμοποιεί διαφορετική, ευρύτερη
μάσκα: η \texttt{\en{cold\_start\_mask}} χαρακτηρίζει κάθε ζεύγος με μηδενικούς
κοινούς γείτονες, ακόμη και όταν και τα δύο άκρα ανήκουν στο γράφημα και το
\en{Preferential Attachment} είναι μη μηδενικό.
Για σαφήνεια το υποσύνολο αυτό
αναφέρεται ως \emph{μηδενικού \en{CN}}.
Το \ColdStartPct\% αφορά αυτό το
διαγνωστικό υποσύνολο και δεν πρέπει να ερμηνεύεται ως ποσοστό ζευγών που
παρακάμπτουν πράγματι το Επίπεδο~1.


\section{Διασταύρωση Επιπέδου και Δυσκολίας}
\label{sec:cascade-diagnostic}

Πέραν της δρομολόγησης ζευγών προς πρόβλεψη, διασταυρώνουμε τα αποτελέσματα του
\en{CascadeLP} με τον χαρακτηρισμό διαχωρισιμότητας του συνόλου δεδομένων, για να
επικυρώσουμε ότι η απόδοσή του είναι γνήσια και όχι τεχνούργημα ρύθμισης ορίων: για κάθε
ζεύγος του συνόλου επαλήθευσης καταγράφουμε ποιο επίπεδο (0--3) το επέλυσε, μαζί
με την ετικέτα δυσκολίας του ζεύγους από τον χαρακτηρισμό διαχωρισιμότητας
(\S\ref{sec:forensics-protocol}). Αυτή η διασταύρωση επιπέδου$\times$δυσκολίας
απαντά σε ένα ερώτημα που η βαθμολογία \en{Macro F1} του \en{CascadeLP} δεν
απαντά από μόνη της: πόσα από τα ζεύγη που φτάνουν στο ακριβό Επίπεδο~3 είναι
πράγματι δύσκολα ζεύγη, και πόσα θα μπορούσαν να είχαν επιλυθεί σε φθηνότερο
επίπεδο με διαφορετική ρύθμιση ορίου;

Για κάθε ζεύγος $(u,v)$ του συνόλου επαλήθευσης καταγράφεται η πλειάδα
$(y_{uv}, \hat{y}_{uv}, t_{uv}, d_{uv})$, όπου $t_{uv} \in \{0,1,2,3\}$ είναι το
επίπεδο που το επέλυσε και $d_{uv}$ είναι η ετικέτα δυσκολίας.
Από αυτές τις
πλειάδες υπολογίζονται δύο πίνακες: η ακρίβεια και η βαθμολογία \en{Macro F1} ανά
επίπεδο, και η ακρίβεια ανά συνδυασμό επιπέδου και δυσκολίας.
Ο δεύτερος πίνακας
είναι ιδιαίτερα κατατοπιστικός: εάν το μεγαλύτερο μέρος των ζευγών στο Επίπεδο~3
είναι ήδη χαρακτηρισμένο ως «δύσκολο» από το πρωτόκολλο ανίχνευσης, τότε η
δρομολόγηση ακολουθεί την αναμενόμενη σειρά των απλών κριτηρίων· εάν αντίθετα
πολλά ζεύγη Επιπέδου~3
είναι χαρακτηρισμένα «τετριμμένα», αυτό υποδηλώνει ότι τα όρια αξιοπιστίας
$\tau_1, \tau_2$ είναι υπερβολικά συντηρητικά.


\section{Μοντέλο Αναφοράς: \en{TF-IDF + SVM}}
\label{sec:svm-baseline}

Ανεξάρτητα από την αλυσίδα \en{CascadeLP}, εκπαιδεύουμε ένα \en{SVM} πυρήνα
ακτινικής βάσης στην ομοιότητα συνημίτονου \en{TF-IDF} μεταξύ των δύο τελικών
κόμβων κάθε ζεύγους, ακολουθώντας το ιστορικό προηγούμενο των
\textcite{alhasan2006link} για επιβλεπόμενη πρόβλεψη ακμών.
Το μοντέλο αυτό δεν
αποτελεί μέρος της αλυσίδας κλιμάκωσης· χρησιμεύει ως ανεξάρτητο σημείο αναφοράς
κλασικής μηχανικής μάθησης για το πόσο απαραίτητα είναι τα εξειδικευμένα μοντέλα
του \en{CascadeLP} έναντι ενός απλούστερου, ενιαίου ταξινομητή.

Το κόστος εκπαίδευσης ενός \en{SVM} με πυρήνα ακτινικής βάσης είναι
$O(n^2)$--$O(n^3)$ ως προς τον αριθμό ζευγών εκπαίδευσης $n$, καθιστώντας την
εκπαίδευση στο πλήρες σύνολο των \TrainPairs{} ζευγών εκπαίδευσης απαγορευτική.
Για
τον λόγο αυτό, το μοντέλο εκπαιδεύεται σε ένα στρωματοποιημένο (κατά ετικέτα)
δείγμα \SvmSampleSize{} ζευγών, μέγεθος αρκετά μεγάλο για σταθερό όριο απόφασης σε έναν
μονοδιάστατο χώρο χαρακτηριστικών, αλλά αρκετά μικρό ώστε η εκπαίδευση να
ολοκληρώνεται σε δευτερόλεπτα.

Σε αντίθεση με τη διακλαδισμένη, πολυεπίπεδη δρομολόγηση του \en{CascadeLP}
(Σχήμα~\ref{fig:cascade}), η διοχέτευση του \en{SVM} είναι μία ενιαία, γραμμική
ακολουθία χωρίς κλιμάκωση ή έξοδο υπό συνθήκη --- κάθε ζεύγος περνά από τα ίδια
τέσσερα βήματα, ανεξάρτητα από το πόσο «εύκολο» ή «δύσκολο» είναι
(Σχήμα~\ref{fig:svm-pipeline}).

\begin{figure}[h!]
    \centering
    \tikzset{
        svmbox/.style={rectangle, draw, rounded corners, minimum width=3.0cm,
            minimum height=1.1cm, align=center},
        svmflow/.style={-{Stealth}, thick},
    }
    \begin{tikzpicture}[node distance=0.9cm and 1.0cm]
        \node[svmbox] (texts) {Κείμενα\\άρθρων $u, v$};
        \node[svmbox, right=of texts] (tfidf) {Διανυσματοποίηση\\\en{TF-IDF}};
        \node[svmbox, right=of tfidf] (cos) {Ομοιότητα\\συνημίτονου};
        \node[svmbox, below=of cos] (sample) {Στρωματοπ.\\υποδείγμα \SvmSampleSize{}};
        \node[svmbox, left=of sample] (svc) {\en{SVC} πυρήνα\\ακτινικής βάσης};
        \node[svmbox, left=of svc] (pred) {Πρόβλεψη\\$\hat{y} \in \{0,1\}$};

        \draw[svmflow] (texts) -- (tfidf);
        \draw[svmflow] (tfidf) -- (cos);
        \draw[svmflow] (cos) -- node[right, font=\small]{εκπαίδευση} (sample);
        \draw[svmflow] (sample) -- (svc);
        \draw[svmflow] (svc) -- (pred);
        \draw[svmflow, dashed] (cos.south) -- ++(0,-0.45) -| ($(svc.north)+(0,0)$)
            node[pos=0.55, above, font=\small, xshift=-1.3cm]{συμπερασμός (όλα τα ζεύγη)};
    \end{tikzpicture}
    \caption{Διοχέτευση του μοντέλου αναφοράς \en{TF-IDF + SVM}: μία γραμμική
    ακολουθία βημάτων, χωρίς κλιμάκωση υπό συνθήκη. Η υποδειγματοληψία
    εφαρμόζεται μόνο κατά την εκπαίδευση (συνεχές βέλος)· κατά τον συμπερασμό
    στο σύνολο επικύρωσης ή ελέγχου το εκπαιδευμένο \en{SVC} εφαρμόζεται σε όλα
    τα ζεύγη (διακεκομμένο βέλος).}
    \label{fig:svm-pipeline}
\end{figure}

Ο Αλγόριθμος~\ref{alg:svm} περιγράφει την εκπαίδευση με υποδειγματοληψία.

\begin{algorithm}[h!]
\caption{Εκπαίδευση \en{SVM} με Στρωματοποιημένη Υποδειγματοληψία}
\label{alg:svm}
\begin{algorithmic}[1]
\Require Ομοιότητες \en{TF-IDF} $s \in \mathbb{R}^n$ και ετικέτες $y$ για $n$
    ζεύγη εκπαίδευσης· μέγεθος υποδείγματος $k=\SvmSampleSize{}$
\Ensure Εκπαιδευμένο \en{SVC} $M$
\If{$n \leq k$}
    \State $(s', y') \gets (s, y)$
\Else
    \State $(s', y') \gets$ στρωματοποιημένο τυχαίο δείγμα $k$ ζευγών από $(s, y)$
        \Comment{κατά ετικέτα, \en{seed}$=$\Seed{}}
\EndIf
\State $s' \gets \text{\en{StandardScaler}}(s')$ \Comment{τυποποίηση μηδενικού μέσου}
\State $M \gets \text{\en{SVC}}(\text{kernel}=\text{\en{rbf}}, C=\SvmC{}, \gamma=\text{\en{\SvmGamma{}}})$
\State $M.\text{fit}(s', y')$
\State \Return $M$
\end{algorithmic}
\end{algorithm}


\section{Πρωτόκολλο Γενίκευσης Μεταξύ Συνόλων Δεδομένων}
\label{sec:generalization-protocol}

Για την πλήρη αξιολόγηση του κατά πόσον οι επιδόσεις στο \en{DSAA 2023}
αντικατοπτρίζουν γνήσια σημασιολογική συλλογιστική ή χαρακτηριστικά ειδικά για
αυτό το σύνολο δεδομένων, σχεδιάζεται ένα πρωτόκολλο γενίκευσης μεταξύ συνόλων
δεδομένων: εκπαίδευση του καλύτερου ταξινομητή κειμένου στο \en{DSAA 2023},
αξιολόγηση χωρίς περαιτέρω εκπαίδευση σε ένα δεύτερο σύνολο δεδομένων \en{TAG}, και
αξιολόγηση μετά από ειδίκευση στο σύνολο εκπαίδευσης του δεύτερου συνόλου
δεδομένων.
Η διαφορά μεταξύ των δύο τελευταίων θα αποτελούσε την ένδειξη
γενίκευσης.

Αυτό το ακριβές πρωτόκολλο μεταφοράς---εκπαίδευση στο ένα σύνολο δεδομένων,
αξιολόγηση χωρίς περαιτέρω εκπαίδευση σε άλλο---\textbf{δεν εκτελείται στην
παρούσα φάση} και παραμένει κατεύθυνση μελλοντικής εργασίας
(Κεφάλαιο~\ref{ch:conclusion}).
Εκτελείται όμως μια στενότερη, αλλά ανεξάρτητα
πληροφοριακή, μορφή επικύρωσης: το ίδιο, αναλλοίωτο \en{CascadeLP}
επανεκπαιδεύεται από την αρχή σε ένα δεύτερο, ανεξάρτητα κατασκευασμένο σύνολο
δεδομένων \en{Wikipedia}, χωρίς το τεχνούργημα δειγματοληψίας που εντοπίστηκε
στο \en{DSAA 2023} (\S\ref{sec:hub-artifact}, Κεφάλαιο~\ref{ch:analysis}).
Αυτό δεν ελέγχει μεταφορά μάθησης μεταξύ συνόλων δεδομένων, αλλά ελέγχει κάτι
εξίσου κρίσιμο για την αξιοπιστία της κύριας συνεισφοράς: αν η αρχιτεκτονική
συμπεριφέρεται όπως σχεδιάστηκε σε ένα σύνολο δεδομένων χωρίς τα ζητήματα
ακεραιότητας του \en{DSAA 2023}.
Τα αποτελέσματα παρουσιάζονται στην
\S\ref{sec:wiki-fresh}.
